\chapter{Conclusioni}
\label{cap:conclusioni}

\section{Risposte alle domande di ricerca}
\label{sec:risposte}

\paragraph{H1 --- Proiezioni \texorpdfstring{$Q,K,V$}{Q,K,V} quantistiche.}
H1 è \nonsupportata. La differenza media di validation loss, quantistico meno
classico, è $+0{,}00543$ (IC al 95\% $[-0{,}01040;+0{,}02126]$, $p=0{,}458$).

\paragraph{H2 --- \emph{Patch embedding} quantistico.}
H2 è \nonsupportata. L'accuracy diminuisce in media di $0{,}02414$ (IC al 95\%
$[-0{,}05432;+0{,}00604]$, $p=0{,}091$). La macro-AUC è inferiore nei cinque seed
($-0{,}01016$, IC $[-0{,}01525;-0{,}00507]$).

\paragraph{H3 --- Regolarizzazione implicita.}
H3 è \nonsupportata. Il delta di gap quantistico meno \code{bounded\_mlp} è
positivo su PathMNIST ($+0{,}148$), BloodMNIST ($+0{,}100$) e DermaMNIST
($+0{,}085$), con segno positivo in 12 delle 15 coppie dataset/seed. Solo
l'intervallo di PathMNIST è interamente positivo.

\paragraph{H4 --- Kernel di fedeltà.}
Per H4 emerge \evidenzacontraria. Rispetto al coseno, la differenza media del
criterio di Fisher è $-0{,}01515$ (IC al 95\% $[-0{,}02379;-0{,}00652]$); rispetto all'RBF
è $-0{,}04168$ (IC $[-0{,}06438;-0{,}01899]$). La validazione SVC sul test
ufficiale produce metriche inferiori per IQP nei tre regimi di campionamento.

\paragraph{Risposta alla domanda generale.} Nessuno dei quattro esperimenti
mostra il vantaggio ipotizzato nelle configurazioni provate. La tesi valuta le
accezioni (b), (c) e (d) della \cref{ssec:quattro-promesse}; non misura uno
speedup computazionale.

\section{Contributi}
\label{sec:contributi-finali}

Il lavoro comprende quattro esperimenti appaiati e il pacchetto
\code{quantum\_framework}, che raccoglie layer PyTorch--PennyLane, caricamento dei
dati, addestramento, diagnostiche e analisi statistica. Configurazioni, seed,
comandi e tabelle complete sono conservati con uno schema dei risultati
versionato. La procedura di riproduzione è nell'appendice~\ref{app:riproduzione}.

\section{Limiti}
\label{sec:limiti}

I circuiti hanno 4--10 qubit, due strati, un solo ansatz e una sola codifica. I
VQC sostituiscono componenti circoscritti di modelli altrimenti classici, su
MedMNIST a $28\times28$ e su un corpus a caratteri. Gli esperimenti 02--03 usano
cinque seed, corrispondenti a un $\dz$ minimo rilevabile di $1{,}68$
(\cref{ssec:potenza}). Gli iperparametri sono stati scelti sulla baseline
classica; negli esperimenti 01--02 viene misurata la composizione
adapter${}+{}$VQC. Tutte le esecuzioni quantistiche sono simulate. Il costo della
simulazione ha determinato la scala dei modelli e dei campioni
(\cref{sec:overhead}).

\section{Sviluppi successivi}
\label{sec:sviluppi}

\begin{enumerate}
  \item Variare la dimensionalità compressa per misurare la dipendenza dal collo
    di bottiglia.
  \item Provare altre codifiche, profondità e ansatz, inclusi \emph{data
    re-uploading} e kernel addestrabili.
  \item Aumentare i seed e definire prima dell'esperimento una soglia di
    equivalenza per un TOST. Con 10 e 20 coppie, il $\dz$ minimo rilevabile
    stimato scende rispettivamente a $1{,}00$ e $0{,}66$.
  \item Selezionare separatamente gli iperparametri delle due varianti con lo
    stesso budget; per l'esperimento~04 questo include la bandwidth dell'RBF.
  \item Eseguire un sottoinsieme su hardware per misurare rumore di
    campionamento, costo dei gradienti e tempi.
\end{enumerate}
